在完成阶段划分与机制总结之后，可以进一步构建一个可解释的预测框架，用于估计未来一段时间内的转型稳定性。

\subsection{框架设计}

本研究不追求高复杂度，而强调解释性。框架输入为曝光能力、角色多样性、平台自主性与韧性恢复四个指标，输出为三类路径判断：稳定扩张、阶段波动与危机后重启。

\subsection{可解释性优势}

与复杂模型相比，规则驱动框架的优势在于变量含义清晰、评分路径可追踪、修改成本低。对于 thesis 示例来说，这种设计尤其适合展示“表述清晰、结构完整、公式与图表并存”的写法。

\subsection{应用示例}

在演示案例中，当平台自主性与持续输出能力同时提高时，模型更倾向于给出“重启成功”的判断；若曝光仍高但自主性不足，则更可能被划入“阶段波动”。
